% --- Archon physics formalization source begin ---
% archon:physics
% archon:covers IPhO2026Problems/problem_IPhO_2026_2_C_4.lean
% archon:source-report reports/ipho_2026/problem_IPhO_2026_2_C_4.source.json
% archon:problem-id IPhO_2026_2
% archon:part-id C.4
% archon:previous-part IPhO_2026_2_C_3
% archon:previous-part-policy natural_language_prerequisite_only

\chapter{Physics problem IPhO\_2026\_2\_C\_4}
\label{ch:IPhO2026Problems_problem_IPhO_2026_2_C_4}

\paragraph{Problem source.}
For the half-cylindrical mirror of radius R, ray A is incident at angle theta
and its reflected line is y = m\_A*x + b\_A.  A neighboring parallel ray B is
incident at theta + Delta theta, with Delta theta much smaller than theta, and
its reflected line is y = m\_B*x + b\_B.  The envelope/intersection of neighboring
rays forms the caustic.  Use Figure 2g and its coordinate convention.

Current subquestion:
For theta << 1, put the caustic in the form Y\_c = v*|X\_c|\textasciicircum{}(p/q) + u. Determine u, v, and the integers p,q.

\paragraph{Current subquestion.}
For theta << 1, put the caustic in the form Y\_c = v*|X\_c|\textasciicircum{}(p/q) + u. Determine u, v, and the integers p,q.

\paragraph{Recorded answer/context.}
u = R/2, v = (3/4)*R\textasciicircum{}(1/3), p = 2, and q = 3.

\paragraph{Figure/image path.}
/root/proposal\_for\_physic/science-mango/ipho\_2026\_source/image/T2\_page-4.png

\paragraph{Reusable previous-part conclusions.}
\begin{itemize}
\item Source C.3. Question: Find the limiting intersection coordinates (X\_c,Y\_c) of the neighboring reflected rays. Reusable conclusions: X\_c = R*sin(theta)\textasciicircum{}3; Y\_c = (R/2)*cos(theta)*(2 - cos(2*theta)). Policy: natural\_language\_prerequisite\_only; do\_not\_import\_Lean\_output
\end{itemize}

\paragraph{Formalization target.}
create a compiling Lean file with sorry bodies at `IPhO2026Problems/problem\_IPhO\_2026\_2\_C\_4.lean`. The Lean declarations must preserve the physical quantities, dimensions or dimensional roles, figure labels, governing-law hypotheses, and final relation expressed by this problem.
Use Mathlib/Physlib names found through LeanExplore where available. If a physics API is missing, introduce faithful local abstractions rather than scalar placeholder aliases.

\begin{definition}[LengthReading]
  \label{decl:physics:IPhO_2026_2_C_4:LengthReading}
  \lean{IPhO2026Problems.IPhO2026_2_C_4.LengthReading}
  A scalar reading of a physical length.
\end{definition}
\begin{proof}
  This is the defining declaration; unfolding it gives the stated typed data or relation.
\end{proof}

\begin{definition}[CubeRootLengthReading]
  \label{decl:physics:IPhO_2026_2_C_4:CubeRootLengthReading}
  \lean{IPhO2026Problems.IPhO2026_2_C_4.CubeRootLengthReading}
  The dimension of the coefficient multiplying a length to the power “2 / 3”.
\end{definition}
\begin{proof}
  This is the defining declaration; unfolding it gives the stated typed data or relation.
\end{proof}

\begin{definition}[ReflectedLineReadout]
  \label{decl:physics:IPhO_2026_2_C_4:ReflectedLineReadout}
  \lean{IPhO2026Problems.IPhO2026_2_C_4.ReflectedLineReadout}
  \uses{decl:physics:IPhO_2026_2_C_4:LengthReading}
  The slope and intercept of a reflected ray in the Cartesian coordinate frame of Figure 2g.
\end{definition}
\begin{proof}
  This is the defining declaration; unfolding it gives the stated typed data or relation.
\end{proof}

\begin{definition}[Figure2gOpticalSystem]
  \label{decl:physics:IPhO_2026_2_C_4:Figure2gOpticalSystem}
  \lean{IPhO2026Problems.IPhO2026_2_C_4.Figure2gOpticalSystem}
  \uses{decl:physics:IPhO_2026_2_C_4:LengthReading, decl:physics:IPhO_2026_2_C_4:ReflectedLineReadout}
  The optical data used in part C.4.
\end{definition}
\begin{proof}
  This is the defining declaration; unfolding it gives the stated typed data or relation.
\end{proof}

\begin{definition}[neighboringIntersectionX]
  \label{decl:physics:IPhO_2026_2_C_4:neighboringIntersectionX}
  \lean{IPhO2026Problems.IPhO2026_2_C_4.neighboringIntersectionX}
  \uses{decl:physics:IPhO_2026_2_C_4:Figure2gOpticalSystem}
  The “x”-coordinate of the intersection of reflected ray “A”, incident at angle “θ”, and neighboring reflected ray “B”, incident at “θ + Δθ”.
\end{definition}
\begin{proof}
  This is the defining declaration; unfolding it gives the stated typed data or relation.
\end{proof}

\begin{definition}[neighboringIntersectionY]
  \label{decl:physics:IPhO_2026_2_C_4:neighboringIntersectionY}
  \lean{IPhO2026Problems.IPhO2026_2_C_4.neighboringIntersectionY}
  \uses{decl:physics:IPhO_2026_2_C_4:Figure2gOpticalSystem, decl:physics:IPhO_2026_2_C_4:neighboringIntersectionX}
  The “y”-coordinate of the same neighboring-ray intersection, obtained from the reflected-line equation “y = m\_A x + b\_A”.
\end{definition}
\begin{proof}
  This is the defining declaration; unfolding it gives the stated typed data or relation.
\end{proof}

\begin{definition}[NeighboringReflectedRaysGenerateCaustic]
  \label{decl:physics:IPhO_2026_2_C_4:NeighboringReflectedRaysGenerateCaustic}
  \lean{IPhO2026Problems.IPhO2026_2_C_4.NeighboringReflectedRaysGenerateCaustic}
  \uses{decl:physics:IPhO_2026_2_C_4:Figure2gOpticalSystem, decl:physics:IPhO_2026_2_C_4:neighboringIntersectionX, decl:physics:IPhO_2026_2_C_4:neighboringIntersectionY}
  The governing caustic-envelope law: the caustic point at angle “θ” is the limit of intersections of reflected rays whose incidence-angle separation “Δθ” tends to zero through nonzero values.
\end{definition}
\begin{proof}
  This is the defining declaration; unfolding it gives the stated typed data or relation.
\end{proof}

\begin{definition}[HasPreviousPartC3Coordinates]
  \label{decl:physics:IPhO_2026_2_C_4:HasPreviousPartC3Coordinates}
  \lean{IPhO2026Problems.IPhO2026_2_C_4.HasPreviousPartC3Coordinates}
  \uses{decl:physics:IPhO_2026_2_C_4:Figure2gOpticalSystem}
  The reusable conclusion of part C.3, stated directly rather than importing that part's Lean output.
\end{definition}
\begin{proof}
  This is the defining declaration; unfolding it gives the stated typed data or relation.
\end{proof}

\begin{theorem}[Physics formalization target]
\label{thm:physics:IPhO_2026_2_C_4:target}
\lean{IPhO2026Problems.IPhO2026_2_C_4.determineSmallAngleCaustic}
\uses{decl:physics:IPhO_2026_2_C_4:LengthReading, decl:physics:IPhO_2026_2_C_4:CubeRootLengthReading, decl:physics:IPhO_2026_2_C_4:ReflectedLineReadout, decl:physics:IPhO_2026_2_C_4:Figure2gOpticalSystem, decl:physics:IPhO_2026_2_C_4:neighboringIntersectionX, decl:physics:IPhO_2026_2_C_4:neighboringIntersectionY, decl:physics:IPhO_2026_2_C_4:NeighboringReflectedRaysGenerateCaustic, decl:physics:IPhO_2026_2_C_4:HasPreviousPartC3Coordinates}
The assigned autoformalize agent should translate this physics problem into Lean declarations in the covered file, with theorem and lemma proof bodies written as `by sorry`.
\end{theorem}
\begin{proof}
This is an autoformalization task, not a proof task. Produce statements that can later be proved by the physics prover without weakening the physical model.
\end{proof}
% --- Archon physics formalization source end ---
